% Compile with LuaLaTeX:
%   lualatex -interaction=nonstopmode -halt-on-error STEP4_GEOMETRY_SCALE.tex
% The page is exactly 16:9. All positions are explicit for stable review.
\documentclass[10pt]{article}
\usepackage[paperwidth=16in,paperheight=9in,margin=0in]{geometry}
\usepackage{fontspec}
\usepackage{amsmath}
\usepackage{tikz}
\usetikzlibrary{arrows.meta,backgrounds,calc,positioning}
\pagestyle{empty}
\setmainfont{Arial}
\hyphenpenalty=10000
\exhyphenpenalty=10000

\definecolor{ink}{HTML}{172033}
\definecolor{muted}{HTML}{526173}
\definecolor{bluefill}{HTML}{EAF3FF}
\definecolor{bluestroke}{HTML}{2F6FB3}
\definecolor{greenfill}{HTML}{EAF8F0}
\definecolor{greenstroke}{HTML}{2E7D5B}
\definecolor{purplefill}{HTML}{F2ECFF}
\definecolor{purplestroke}{HTML}{6842A8}
\definecolor{amberfill}{HTML}{FFF5DC}
\definecolor{amberstroke}{HTML}{A66B10}
\definecolor{orangefill}{HTML}{FFF0E6}
\definecolor{orangestroke}{HTML}{B85C2E}
\definecolor{redfill}{HTML}{FDECEC}
\definecolor{redstroke}{HTML}{A34848}
\definecolor{slatefill}{HTML}{F2F5F8}
\definecolor{slatestroke}{HTML}{657487}

\newcommand{\twobullets}[2]{%
  \parbox{2.12in}{\raggedright\fontsize{12}{14}\selectfont
  \textbullet\ #1\\\textbullet\ #2}}
\newcommand{\threebullets}[3]{%
  \parbox{2.12in}{\raggedright\fontsize{12}{14}\selectfont
  \textbullet\ #1\\\textbullet\ #2\\\textbullet\ #3}}
\newcommand{\inputcontent}[2]{%
  \parbox{2.18in}{\raggedright
  \textbf{\fontsize{12.5}{14.5}\selectfont #1}\\[3pt]
  {\fontsize{11.5}{13.2}\selectfont #2}}}

\begin{document}
\begin{tikzpicture}[
  remember picture,
  overlay,
  shift={(current page.south west)},
  x=1in,
  y=1in,
  every node/.style={font=\fontsize{12}{14}\selectfont,text=ink,align=center},
  stage/.style={
    draw,
    rounded corners=7pt,
    line width=1.25pt,
    minimum height=1.52in,
    text width=2.12in,
    inner xsep=8pt,
    inner ysep=9pt,
    fill=greenfill,
    draw=greenstroke
  },
  input/.style={
    draw,
    rounded corners=5pt,
    line width=1.1pt,
    minimum height=0.92in,
    text width=2.18in,
    inner xsep=6pt,
    inner ysep=6pt,
    align=left
  },
  output/.style={
    draw,
    rounded corners=7pt,
    line width=1.35pt,
    minimum height=1.65in,
    text width=2.12in,
    inner xsep=8pt,
    inner ysep=9pt,
    fill=purplefill,
    draw=purplestroke
  },
  flow/.style={
    -{Latex[length=3.5mm,width=2.2mm]},
    line width=1.45pt,
    draw=ink,
    rounded corners=8pt
  },
  reference/.style={
    -{Latex[length=3.1mm,width=2.0mm]},
    line width=1.3pt,
    draw=purplestroke,
    densely dashed,
    rounded corners=7pt
  },
  edgelabel/.style={
    fill=none,
    inner sep=1pt,
    font=\fontsize{14}{16}\selectfont,
    text=ink
  },
  smalllabel/.style={
    fill=none,
    inner sep=1pt,
    font=\fontsize{11.5}{13}\selectfont,
    text=muted
  }
]

\path[use as bounding box] (0,0) rectangle (16,9);

% Header
\node[font=\bfseries\fontsize{20}{23}\selectfont,anchor=north west]
  at (0.40,8.78) {Step 4 - Geometry and Metric-Scale Hypotheses};
\node[font=\fontsize{11.5}{13}\selectfont,text=muted,anchor=north west]
  at (0.42,8.40) {Transform track-aligned image evidence into validated camera and object geometry};
\node[font=\bfseries\fontsize{11}{13}\selectfont,text=purplestroke,anchor=north east]
  at (15.60,8.69) {Output: ranked geometry hypotheses with uncertainty};

% Input panel
\begin{scope}[on background layer]
  \filldraw[fill=bluefill!34,draw=bluestroke!65,line width=1.2pt,rounded corners=10pt]
    (0.38,1.47) rectangle (3.25,8.02);
\end{scope}
\node[font=\bfseries\fontsize{14}{16}\selectfont,anchor=west]
  at (0.62,7.78) {Inputs};
\node[font=\fontsize{12}{14}\selectfont,text=muted,anchor=west]
  at (0.62,7.45) {$\mathcal{I}_4=\{\mathcal{T},r_x,\mathcal{K}\}$};

\node[input,fill=purplefill,draw=purplestroke] (package) at (1.81,6.57)
  {\inputcontent{Tracking Package $\mathcal{T}$}
  {track view, mask candidates, state and association audit}};

\node[input,minimum height=1.18in,fill=bluefill,draw=bluestroke] (store) at (1.81,5.13)
  {\inputcontent{Evidence Store $r_x$}
  {RGB, masks/logits, flow, depth, confidence and transforms}};

\node[input,fill=amberfill,draw=amberstroke] (priors) at (1.81,3.59)
  {\inputcontent{Metadata + Priors $\mathcal{K}$}
  {camera data when available; frozen height, size and road ranges}};

\node[draw=slatestroke,fill=slatefill,rounded corners=5pt,line width=1.0pt,
      text width=2.18in,minimum height=0.86in,inner sep=6pt]
  (gate) at (1.81,2.15)
  {\textbf{Retention Gate}\\[3pt]
   {\fontsize{11.5}{13}\selectfont require references, transforms and hashes to resolve}};

\draw[reference] (package.south) -- node[smalllabel,right=6pt,text=purplestroke]
  {$r_x$} (store.north);

% Processing graph
\node[stage] (align) at (4.75,5.83)
  {\textbf{\fontsize{13.5}{15.5}\selectfont 4.1 - Resolve + Align Evidence}\\[7pt]
   \threebullets{resolve immutable artifact refs}
                {align time and coordinates}
                {preserve candidates + uncertainty}};

\node[stage,fill=greenfill!72] (pose) at (7.55,6.86)
  {\textbf{\fontsize{13.5}{15.5}\selectfont 4.2 - Camera Motion}\\[7pt]
   \threebullets{exclude tracked foreground}
                {fit motion on background}
                {retain pose candidates + covariance}};

\node[stage,fill=bluefill,draw=bluestroke] (ground) at (7.55,4.55)
  {\textbf{\fontsize{13.5}{15.5}\selectfont 4.3 - Intrinsics + Ground}\\[7pt]
   \threebullets{validate or estimate $K$}
                {fit horizon / road plane}
                {extract ground-contact anchors}};

\node[stage,minimum height=1.72in,fill=amberfill,draw=amberstroke] (scale) at (10.47,5.83)
  {\textbf{\fontsize{13.5}{15.5}\selectfont 4.4 - Scale Hypotheses}\\[7pt]
   \threebullets{use metric depth only if validated}
                {camera-height / size / road priors}
                {emit scale status + interval}};

\node[stage,minimum height=1.72in,fill=orangefill,draw=orangestroke] (lift) at (13.35,5.83)
  {\textbf{\fontsize{13.5}{15.5}\selectfont 4.5 - Lift to Candidate 3D}\\[7pt]
   \twobullets{consolidate masked depth + anchors}
              {back-project and transform with covariance}\\[6pt]
   {\fontsize{11.5}{13}\selectfont
    $X_c=zK^{-1}p,\quad X_w=R_{wc}X_c+t_{wc}$}};

\node[stage,text width=2.32in,minimum height=1.70in,fill=redfill,draw=redstroke]
  (validate) at (10.47,2.95)
  {\textbf{\fontsize{13.5}{15.5}\selectfont 4.6 - Geometric Validation}\\[7pt]
   \threebullets{mask / box reprojection residuals}
                {flow, depth and ground agreement}
                {reject hard failures; rank survivors}};

\node[output] (geometry) at (13.46,2.95)
  {\textbf{\fontsize{13.5}{15.5}\selectfont 4.7 - Geometry Hypothesis Set}\\[7pt]
   \threebullets{camera pose + scale interval}
                {per-track 3D observations}
                {covariance, residuals + provenance}};

% Input bundle. Raw arrays are read through refs rather than copied into tracks.
\coordinate (inputout) at (3.25,5.83);
\draw[flow] (inputout) -- node[edgelabel,above=7pt] {$\mathcal{I}_4$} (align.west);

% Alignment splits into two complementary estimators.
\draw[flow] (align.east) .. controls +(0.36,0.16) and +(-0.42,-0.06) ..
  node[smalllabel,pos=0.58,above=7pt] {$E^{\mathrm{bg}}$} (pose.west);
\draw[flow] (align.east) .. controls +(0.36,-0.16) and +(-0.42,0.06) ..
  node[smalllabel,pos=0.58,below=7pt] {$E^{\mathrm{geo}}$} (ground.west);

% Pose and ground/calibration candidates jointly constrain scale.
\draw[flow] (pose.east) .. controls +(0.42,0.00) and +(-0.42,0.30) ..
  node[edgelabel,pos=0.52,above=8pt] {$\mathcal{P}$} (scale.north west);
\draw[flow] (ground.east) .. controls +(0.42,0.00) and +(-0.42,-0.30) ..
  node[edgelabel,pos=0.52,below=8pt] {$\mathcal{C}$} (scale.south west);

\draw[flow] (scale.east) -- node[edgelabel,above=8pt] {$s^{1:k}$} (lift.west);

% 3D candidates descend into validation and then become the Step 4 output.
\draw[flow] (lift.south) -- ++(0,-0.43) -| node[edgelabel,pos=0.62,above=9pt]
  {$\mathcal{X}^{1:k}$} (validate.north);
\draw[flow] (validate.east) -- node[edgelabel,above=8pt] {$\mathcal{G}^{1:k}$} (geometry.west);

\coordinate (stepfive) at (15.80,2.95);
\draw[flow] (geometry.east) -- node[edgelabel,above=7pt] {$\mathcal{G}$} (stepfive);

% Boundary note
\node[draw=redstroke,fill=redfill,rounded corners=5pt,line width=1.0pt,
      text width=9.60in,minimum height=0.66in,inner sep=6pt,
      font=\bfseries\fontsize{11.5}{13.5}\selectfont]
  (boundary) at (9.48,1.35)
  {Observability boundary: monocular scale may be relative, ambiguous, or unobservable. Report a hypothesis set and interval; claim metric 3D only when scale is supported and validated.};

\node[anchor=south west,font=\fontsize{11.5}{13}\selectfont,text=muted]
  at (3.90,0.68) {$E^{\mathrm{bg}}$ background evidence · $E^{\mathrm{geo}}$ geometry evidence ·
  $\mathcal{P}$ pose · $\mathcal{C}$ calibration/ground · $s$ scale · $\mathcal{X}$ 3D candidates · $\mathcal{G}$ geometry};

% Footer
\node[anchor=south west,font=\fontsize{10.5}{12.5}\selectfont,text=muted]
  at (0.42,0.30) {Step 4 reads Step 3 artifacts by reference; it does not rerun YOLO, SAM 2, RAFT or DA3.};
\node[anchor=south east,font=\fontsize{10.5}{12.5}\selectfont,text=muted]
  at (15.58,0.30) {Metric motion states and physical trajectories are estimated in Step 5.};

\end{tikzpicture}
\end{document}
